% to be included in the preamble
% use % to be included in the preamble
% use % to be included in the preamble
% use % to be included in the preamble
% use \input{math_preamble} in the preamble of the main document

% math & font packages
\usepackage{amssymb}
\usepackage{mathtools} % \DeclarePairedDelimiter
\usepackage{bm} % bold greek

% bold vectors for each alphabet vx
\usepackage{forloop}
\newcommand{\defvec}[1]{\expandafter\newcommand\csname v#1\endcsname{{\mathbf{#1}}}}
\newcounter{ct}
\forLoop{1}{26}{ct}{
    \edef\letter{\alph{ct}}
    \expandafter\defvec\letter
}

% captial \vA
\forLoop{1}{26}{ct}{
    \edef\letter{\Alph{ct}}
    \expandafter\defvec\letter
}

% math
\DeclareMathOperator{\diag}{diag}
\DeclareMathOperator*{\tr}{tr} % trace
\DeclareMathOperator*{\var}{var}
\DeclareMathOperator*{\cov}{cov}

\DeclarePairedDelimiter{\abs}{\lvert}{\rvert}
\DeclarePairedDelimiter{\norm}{\lVert}{\rVert}

\newcommand{\dm}[1]{\ensuremath{\mathrm{d}{#1}}} % dx dy dz dmu
%\newcommand{\RN}[2]{\frac{\dm{#1}}{\dm{#2}}} % (Radon-Nikodym) derivative
\newcommand{\PD}[2]{\frac{\partial #1}{\partial #2}} % partial derivative

\newcommand{\field}[1]{\ensuremath{\mathbb{#1}}}
\newcommand{\reals}{\field{R}}
\newcommand{\complex}{\field{C}}
\newcommand{\integers}{\field{Z}}
\newcommand{\naturalNumbers}{\field{N}}
\newcommand{\trp}{{^\top}} % transpose
\newcommand{\identity}{\ensuremath{\mathbb{I}}}
\newcommand{\ones}{\ensuremath{\mathbf{1}}}
\newcommand{\indicator}[1]{\mathbf{1}_{#1}} % indicator function
\newcommand{\stateSpace}{\reals^n}
\newcommand{\DF}{\nabla_{\vx}\vf} %\newcommand{\DF}{\mathcal{D}\vf}
\newcommand{\vDelta}{\mathbf{\Delta}}
\newcommand{\vPsi}{\mathbf{\Psi}}
\newcommand{\vphi}{\bm{\phi}}

% Exercises and solutions: shared machinery (toggle + printer-friendly boxes).
\input{exercise_setup}
 in the preamble of the main document

% math & font packages
\usepackage{amssymb}
\usepackage{mathtools} % \DeclarePairedDelimiter
\usepackage{bm} % bold greek

% bold vectors for each alphabet vx
\usepackage{forloop}
\newcommand{\defvec}[1]{\expandafter\newcommand\csname v#1\endcsname{{\mathbf{#1}}}}
\newcounter{ct}
\forLoop{1}{26}{ct}{
    \edef\letter{\alph{ct}}
    \expandafter\defvec\letter
}

% captial \vA
\forLoop{1}{26}{ct}{
    \edef\letter{\Alph{ct}}
    \expandafter\defvec\letter
}

% math
\DeclareMathOperator{\diag}{diag}
\DeclareMathOperator*{\tr}{tr} % trace
\DeclareMathOperator*{\var}{var}
\DeclareMathOperator*{\cov}{cov}

\DeclarePairedDelimiter{\abs}{\lvert}{\rvert}
\DeclarePairedDelimiter{\norm}{\lVert}{\rVert}

\newcommand{\dm}[1]{\ensuremath{\mathrm{d}{#1}}} % dx dy dz dmu
%\newcommand{\RN}[2]{\frac{\dm{#1}}{\dm{#2}}} % (Radon-Nikodym) derivative
\newcommand{\PD}[2]{\frac{\partial #1}{\partial #2}} % partial derivative

\newcommand{\field}[1]{\ensuremath{\mathbb{#1}}}
\newcommand{\reals}{\field{R}}
\newcommand{\complex}{\field{C}}
\newcommand{\integers}{\field{Z}}
\newcommand{\naturalNumbers}{\field{N}}
\newcommand{\trp}{{^\top}} % transpose
\newcommand{\identity}{\ensuremath{\mathbb{I}}}
\newcommand{\ones}{\ensuremath{\mathbf{1}}}
\newcommand{\indicator}[1]{\mathbf{1}_{#1}} % indicator function
\newcommand{\stateSpace}{\reals^n}
\newcommand{\DF}{\nabla_{\vx}\vf} %\newcommand{\DF}{\mathcal{D}\vf}
\newcommand{\vDelta}{\mathbf{\Delta}}
\newcommand{\vPsi}{\mathbf{\Psi}}
\newcommand{\vphi}{\bm{\phi}}

% Exercises and solutions: shared machinery (toggle + printer-friendly boxes).
% Shared exercise/solution machinery, \input by both latent_dynamics_notes.tex
% (via math_preamble.tex) and kalman_filter_exercise.tex. Kept in one file so
% the toggle and box styling are not duplicated.
%
% Single source, student/instructor toggle:
%   Student PDF (default): solution boxes are dropped entirely.
%   Instructor PDF: built with \def\SolutionsBuild{} prepended (see the
%   Makefile *_solutions targets). Do NOT hand-edit the toggle here.
% Boxes are printer-friendly: light tint fills (<= 12%) that stay legible in
% grayscale. Colors are defined here so the file is self-contained.

\usepackage{environ}   % capture-or-discard env bodies (robust inside lists)
\usepackage[most]{tcolorbox}
\definecolor{solnframe}{rgb}{0.10,0.25,0.55} % calm blue for solution boxes
\newif\ifsolutions
\ifdefined\SolutionsBuild\solutionstrue\fi

% Inline micro-exercise, numbered per section, shown in both builds. Optional
% argument tags it, e.g. \begin{exercise}[optional stretch] for the harder tier.
\newtcolorbox[auto counter, number within=section]{exercise}[1][]{%
    breakable, enhanced, sharp corners,
    colback=black!3, colframe=black!55, boxrule=0.5pt,
    left=6pt, right=6pt, top=3pt, bottom=3pt,
    coltitle=black, colbacktitle=black!10, fonttitle=\bfseries,
    title={Exercise~\thetcbcounter\if\relax\detokenize{#1}\relax\else\ (#1)\fi},
}

% Solution box, instructor build only. In the student build the environment
% body is captured by environ and simply not emitted (discarded, not scanned),
% which is robust even inside enumerate/itemize.
\newtcolorbox{solnbox}{%
    breakable, enhanced, sharp corners,
    colback=solnframe!5, colframe=solnframe!70, boxrule=0.5pt,
    left=6pt, right=6pt, top=3pt, bottom=3pt,
    coltitle=solnframe!85!black, colbacktitle=solnframe!12,
    fonttitle=\bfseries, title={Solution},
}
\ifsolutions
    \NewEnviron{solution}{\begin{solnbox}\BODY\end{solnbox}}
\else
    \NewEnviron{solution}{}% discard body in the student build
\fi

 in the preamble of the main document

% math & font packages
\usepackage{amssymb}
\usepackage{mathtools} % \DeclarePairedDelimiter
\usepackage{bm} % bold greek

% bold vectors for each alphabet vx
\usepackage{forloop}
\newcommand{\defvec}[1]{\expandafter\newcommand\csname v#1\endcsname{{\mathbf{#1}}}}
\newcounter{ct}
\forLoop{1}{26}{ct}{
    \edef\letter{\alph{ct}}
    \expandafter\defvec\letter
}

% captial \vA
\forLoop{1}{26}{ct}{
    \edef\letter{\Alph{ct}}
    \expandafter\defvec\letter
}

% math
\DeclareMathOperator{\diag}{diag}
\DeclareMathOperator*{\tr}{tr} % trace
\DeclareMathOperator*{\var}{var}
\DeclareMathOperator*{\cov}{cov}

\DeclarePairedDelimiter{\abs}{\lvert}{\rvert}
\DeclarePairedDelimiter{\norm}{\lVert}{\rVert}

\newcommand{\dm}[1]{\ensuremath{\mathrm{d}{#1}}} % dx dy dz dmu
%\newcommand{\RN}[2]{\frac{\dm{#1}}{\dm{#2}}} % (Radon-Nikodym) derivative
\newcommand{\PD}[2]{\frac{\partial #1}{\partial #2}} % partial derivative

\newcommand{\field}[1]{\ensuremath{\mathbb{#1}}}
\newcommand{\reals}{\field{R}}
\newcommand{\complex}{\field{C}}
\newcommand{\integers}{\field{Z}}
\newcommand{\naturalNumbers}{\field{N}}
\newcommand{\trp}{{^\top}} % transpose
\newcommand{\identity}{\ensuremath{\mathbb{I}}}
\newcommand{\ones}{\ensuremath{\mathbf{1}}}
\newcommand{\indicator}[1]{\mathbf{1}_{#1}} % indicator function
\newcommand{\stateSpace}{\reals^n}
\newcommand{\DF}{\nabla_{\vx}\vf} %\newcommand{\DF}{\mathcal{D}\vf}
\newcommand{\vDelta}{\mathbf{\Delta}}
\newcommand{\vPsi}{\mathbf{\Psi}}
\newcommand{\vphi}{\bm{\phi}}

% Exercises and solutions: shared machinery (toggle + printer-friendly boxes).
% Shared exercise/solution machinery, \input by both latent_dynamics_notes.tex
% (via math_preamble.tex) and kalman_filter_exercise.tex. Kept in one file so
% the toggle and box styling are not duplicated.
%
% Single source, student/instructor toggle:
%   Student PDF (default): solution boxes are dropped entirely.
%   Instructor PDF: built with \def\SolutionsBuild{} prepended (see the
%   Makefile *_solutions targets). Do NOT hand-edit the toggle here.
% Boxes are printer-friendly: light tint fills (<= 12%) that stay legible in
% grayscale. Colors are defined here so the file is self-contained.

\usepackage{environ}   % capture-or-discard env bodies (robust inside lists)
\usepackage[most]{tcolorbox}
\definecolor{solnframe}{rgb}{0.10,0.25,0.55} % calm blue for solution boxes
\newif\ifsolutions
\ifdefined\SolutionsBuild\solutionstrue\fi

% Inline micro-exercise, numbered per section, shown in both builds. Optional
% argument tags it, e.g. \begin{exercise}[optional stretch] for the harder tier.
\newtcolorbox[auto counter, number within=section]{exercise}[1][]{%
    breakable, enhanced, sharp corners,
    colback=black!3, colframe=black!55, boxrule=0.5pt,
    left=6pt, right=6pt, top=3pt, bottom=3pt,
    coltitle=black, colbacktitle=black!10, fonttitle=\bfseries,
    title={Exercise~\thetcbcounter\if\relax\detokenize{#1}\relax\else\ (#1)\fi},
}

% Solution box, instructor build only. In the student build the environment
% body is captured by environ and simply not emitted (discarded, not scanned),
% which is robust even inside enumerate/itemize.
\newtcolorbox{solnbox}{%
    breakable, enhanced, sharp corners,
    colback=solnframe!5, colframe=solnframe!70, boxrule=0.5pt,
    left=6pt, right=6pt, top=3pt, bottom=3pt,
    coltitle=solnframe!85!black, colbacktitle=solnframe!12,
    fonttitle=\bfseries, title={Solution},
}
\ifsolutions
    \NewEnviron{solution}{\begin{solnbox}\BODY\end{solnbox}}
\else
    \NewEnviron{solution}{}% discard body in the student build
\fi

 in the preamble of the main document

% math & font packages
\usepackage{amssymb}
\usepackage{mathtools} % \DeclarePairedDelimiter
\usepackage{bm} % bold greek

% bold vectors for each alphabet vx
\usepackage{forloop}
\newcommand{\defvec}[1]{\expandafter\newcommand\csname v#1\endcsname{{\mathbf{#1}}}}
\newcounter{ct}
\forLoop{1}{26}{ct}{
    \edef\letter{\alph{ct}}
    \expandafter\defvec\letter
}

% captial \vA
\forLoop{1}{26}{ct}{
    \edef\letter{\Alph{ct}}
    \expandafter\defvec\letter
}

% math
\DeclareMathOperator{\diag}{diag}
\DeclareMathOperator*{\tr}{tr} % trace
\DeclareMathOperator*{\var}{var}
\DeclareMathOperator*{\cov}{cov}

\DeclarePairedDelimiter{\abs}{\lvert}{\rvert}
\DeclarePairedDelimiter{\norm}{\lVert}{\rVert}

\newcommand{\dm}[1]{\ensuremath{\mathrm{d}{#1}}} % dx dy dz dmu
%\newcommand{\RN}[2]{\frac{\dm{#1}}{\dm{#2}}} % (Radon-Nikodym) derivative
\newcommand{\PD}[2]{\frac{\partial #1}{\partial #2}} % partial derivative

\newcommand{\field}[1]{\ensuremath{\mathbb{#1}}}
\newcommand{\reals}{\field{R}}
\newcommand{\complex}{\field{C}}
\newcommand{\integers}{\field{Z}}
\newcommand{\naturalNumbers}{\field{N}}
\newcommand{\trp}{{^\top}} % transpose
\newcommand{\identity}{\ensuremath{\mathbb{I}}}
\newcommand{\ones}{\ensuremath{\mathbf{1}}}
\newcommand{\indicator}[1]{\mathbf{1}_{#1}} % indicator function
\newcommand{\stateSpace}{\reals^n}
\newcommand{\DF}{\nabla_{\vx}\vf} %\newcommand{\DF}{\mathcal{D}\vf}
\newcommand{\vDelta}{\mathbf{\Delta}}
\newcommand{\vPsi}{\mathbf{\Psi}}
\newcommand{\vphi}{\bm{\phi}}

% Exercises and solutions: shared machinery (toggle + printer-friendly boxes).
% Shared exercise/solution machinery, \input by both latent_dynamics_notes.tex
% (via math_preamble.tex) and kalman_filter_exercise.tex. Kept in one file so
% the toggle and box styling are not duplicated.
%
% Single source, student/instructor toggle:
%   Student PDF (default): solution boxes are dropped entirely.
%   Instructor PDF: built with \def\SolutionsBuild{} prepended (see the
%   Makefile *_solutions targets). Do NOT hand-edit the toggle here.
% Boxes are printer-friendly: light tint fills (<= 12%) that stay legible in
% grayscale. Colors are defined here so the file is self-contained.

\usepackage{environ}   % capture-or-discard env bodies (robust inside lists)
\usepackage[most]{tcolorbox}
\definecolor{solnframe}{rgb}{0.10,0.25,0.55} % calm blue for solution boxes
\newif\ifsolutions
\ifdefined\SolutionsBuild\solutionstrue\fi

% Inline micro-exercise, numbered per section, shown in both builds. Optional
% argument tags it, e.g. \begin{exercise}[optional stretch] for the harder tier.
\newtcolorbox[auto counter, number within=section]{exercise}[1][]{%
    breakable, enhanced, sharp corners,
    colback=black!3, colframe=black!55, boxrule=0.5pt,
    left=6pt, right=6pt, top=3pt, bottom=3pt,
    coltitle=black, colbacktitle=black!10, fonttitle=\bfseries,
    title={Exercise~\thetcbcounter\if\relax\detokenize{#1}\relax\else\ (#1)\fi},
}

% Solution box, instructor build only. In the student build the environment
% body is captured by environ and simply not emitted (discarded, not scanned),
% which is robust even inside enumerate/itemize.
\newtcolorbox{solnbox}{%
    breakable, enhanced, sharp corners,
    colback=solnframe!5, colframe=solnframe!70, boxrule=0.5pt,
    left=6pt, right=6pt, top=3pt, bottom=3pt,
    coltitle=solnframe!85!black, colbacktitle=solnframe!12,
    fonttitle=\bfseries, title={Solution},
}
\ifsolutions
    \NewEnviron{solution}{\begin{solnbox}\BODY\end{solnbox}}
\else
    \NewEnviron{solution}{}% discard body in the student build
\fi

